\documentclass{article}
\usepackage{graphicx} % Required for inserting images
\usepackage{geometry}
\usepackage{listings}
\usepackage{xcolor}
\usepackage{minted}
\usepackage{amsmath}
\title{llmnote}
\author{jhlyu}
\date{\today}
\newcommand{\bs}{\textrm{batch\_size}}
\newcommand{\seq}{\textrm{seq}}
\newcommand{\dm}{\textrm{d\_model}}
\newcommand{\vs}{\textrm{vocab\_size}}
\newcommand{\fea}{\textrm{feature}}
\newcommand{\raw}{\textrm{raw}}


\begin{document}

\maketitle

\section{Introduction}
Prediction, prediction, prediction! Let us consider a general model for prediction. We have a seq \(x_{0}x_{1}\dots x_{n-1}\) and we want to predict the next symbol $x_{n}$. An easy example would be we see a all 0 seq \(x_{0}x_{1}\dots x_{n-1}=00\dots0\). We naturally predict that the next digit \(x_{n}\) is 0 with large confidence. Mathematically speaking, a prediction machine provides the probability of the next symbol(token) based on all previous token:
\begin{align}
    \text{a prediction machine} \sim \Pr(x_{n}|x_{0}\dots x_{n-1})~.
\end{align}
The main question is how to construct this conditional probability function. Or we want to build a raw logits function of \(x_{n}\) = \(z(x_{0}x_{1}\dots x_{n-1})\). A first straightforward step is to embed these tokens into high dimension space. But what is the next? It will be super inefficient to train a series of function \(\{z_i( x_0 \dots x_{i-1})\}_{i=1}\) to predict the \(i\)-th token \(X_{i}\). A different perspective is to use correlation function to predict the next token. More precisely, we will not use the regular correlation but a pairwise similarity function which can be viewed as correlation-like. I think it would be better to start from a specific example, a sentence:

\begin{align}
    \boxed{\text{The cat sat on the mat because } \textcolor{red}{\text{it}} \text{ was tired.}}~.
\end{align}
The red "it" refers to "the cat". But how should we design a mechanism to let the machine realize it?

\section{Tech Details}
In this section, we review the technical part of the transformer architecture. We mainly focus on the self-attention part. The basic Q is to how do we predict the next token based on the current information. Let us say that the data we have is $X^{\raw} \in (\bs \times \seq ) $. In each entries we have an ID. The first thing we do is to embedding each ID into a higher dimension $\dm$. The matrix is denoted as $E \in (\vs \times \dm)$. The embedding is simply
\begin{minted}[
    linenos,
    frame=single,
    breaklines,
    fontsize=\small
]{python}
class embedding(nn.Module):
    def __init__(self, vocab_size, d_model, weights):
        super().__init__()
        self.weights = nn.Parameter(weights)
    def forward(self, token_ids):
        return self.weights[token_ids]
\end{minted}
Now we have input feature $X \in ( \bs \times\seq\times\dm)$. To simplify the notation without losing any generality, we set $\bs = 1$, i.e., 
\begin{align}
    X = \begin{pmatrix}
        e_1 \\
        e_2 \\
        \vdots\\
        e_n-1
    \end{pmatrix}~.
\end{align}
The attention mechanism requires three components, $W_Q,W_K,W_V$. Here, the big picture is that $W_Q,W_K,W_V$ matrices assign three vectors $q=W_{Q}^{\top} x,  k=W_{K}^{\top}x , v=W_{V}^{\top} x$ to each embedding vector $x$. We list the shape of these three matrices:
\begin{align*}
    &W_Q \in (\dm, \text{d}_k) \\
    &W_K \in (\dm, \text{d}_k) \\
    &W_V \in (\dm, \text{d}_v) ~.
\end{align*}
$W_Q, W_K$ map each token $x$ to vectors $q_{x},k_{x}$ living in $d_{k}$ dimension, respectively. We use these two sets of vector to compute how tokens are connected. We build a score product table:
\begin{align}
    (\text{score})_{ij} = \frac{1}{\sqrt{d_{k}}} q_{x_i}\cdot k_{x_{j}}~.
\end{align}
This can show the similarity among tokens. Vaguely speaking, the large entry \((\text{score})_{ij}\) in this table means that the token \(i\) and the token \(j\) are strong connected. Then we do a softmax along the last dimension of which results are called attention:
\begin{align}
    (\textrm{attn})_{ij} = \exp[ (\textrm{score})_{ij}]/\sum_{k}\exp[ (\textrm{score})_{ik}]~.
\end{align}
For each $(\textrm{attn})_{ij}$, it has the property of
\begin{align}
    \sum_{j}(\textrm{attn})_{ij} = 1~.
\end{align}
To do real prediction, we need to mask out each $(\textrm{score})_{ij}$ with $i > j$. The reason is that for each query \(i\) we can only look up the key backward. So we need set all $(\textrm{score})_{ij}$ with $i > j$ to be $-\infty$. For example, a proper attn matrix looks like 
\[
A=
\begin{array}{c|cccccccccc}
 & \text{The} & \text{cat} & \text{sat} & \text{on} & \text{the}
 & \text{mat} & \text{because} & \text{it} & \text{was} & \text{tired}\\
\hline

\text{The}
&1.000&0&0&0&0&0&0&0&0&0\\

\text{cat}
&0.366&0.634&0&0&0&0&0&0&0&0\\

\text{sat}
&0.131&0.413&0.456&0&0&0&0&0&0&0\\

\text{on}
&0.120&0.163&0.296&0.421&0&0&0&0&0&0\\

\text{the}
&0.227&0.113&0.107&0.196&0.357&0&0&0&0&0\\

\text{mat}
&0.079&0.101&0.130&0.184&0.203&0.303&0&0&0&0\\

\text{because}
&0.081&0.115&0.155&0.099&0.089&0.133&0.328&0&0&0\\

\text{it}
&0.059&0.290&0.097&0.053&0.056&0.124&0.107&0.215&0&0\\

\text{was}
&0.056&0.097&0.088&0.062&0.059&0.072&0.097&0.217&0.252&0\\

\text{tired}
&0.039&0.113&0.065&0.044&0.041&0.062&0.072&0.177&0.160&0.227
\end{array}~.
\]
This can be viewed as correlation information contained in this sentence. We do not want to sent this whole attn matrix into our machine learning model. Instead, we introduce a mechanism to propagate the information forward such that we only predict the next token based on the modified last token \(\Tilde{x}_{n-1}\).
\begin{align}
    &z(x_0x_1\dots x_{n-1}) \leftarrow \text{not the function we want learn} \\
    &z(\Tilde{x}_{n-1})  \leftarrow \text{the function we want learn} ~.
\end{align}
To do this, we decorate the original feature $X$. We introduce another vector \(x_v \in \text{d}_v\) for each token. For the whole feature, the $V$ matrix is $V=XW_V \in (\seq \times \text{d}_v)$. The attn is simply \( \text{attn} = A@V\)
\begin{align}
\mathrm{Attn}
&=
A V \\
&=
\begin{pmatrix}
1.000\,v_1 \\
0.366\,v_1 + 0.634\,v_2 \\
0.131\,v_1 + 0.413\,v_2 + 0.456\,v_3 \\
0.120\,v_1 + 0.163\,v_2 + 0.296\,v_3 + 0.421\,v_4 \\
\vdots
\end{pmatrix}.
\end{align}
This is for single head attention $\textrm{attn}\in (\seq, d_{v})$. In piratical, we can have multiple heads.
\begin{align}
    (\textrm{head}_1, \dots,\textrm{head}_{h}) \in (\bs, \seq, d_{v}\times h= \dm)~. 
\end{align}
Then we mix this multi-head attention with a matrix $O\in(\dm, \dm)$. And add this back to $X$
\begin{align}
    X' = X + \textrm{Multi-Head}(X)~.
\end{align} 
And then add FNN to this X':
\begin{align}
    X'' = X' + \textrm{FNN}(X')
\end{align}
Here we list code.
\end{document}